\subsection{Feast: reconstructing the information available to a banking decision}
\label{exec19:feast-clock}

Consider a credit team replaying yesterday's offer decisions. A transaction
can have an early business date, a later correction date and a still later
arrival in the system used by the offer engine. Joining on the business date
alone does not reproduce the information on which the original offer could
depend. Feast is a concrete candidate for organizing feature definitions and
retrieving historical and online features, but its usefulness for this
proposal depends on a precise observation contract. The selected original
revision is commit \texttt{9d42729d471a} of 7 October 2026. Its historical API
already provides an optional creation-time cutoff; the appropriate comparison
must acknowledge that capability and investigate its meaning for each
provider.\citep{exec19_feast_source}

For a feature version $j$ associated with customer $i$, let $e_{ij}$ be its
event time, $c_{ij}$ its source creation time and $a_{ij}$ the first time that
the version, with the required transformations, was usable by the decision
system. Let $t_i$ denote the decision time and $\ell>0$ a declared feature
lifetime. The proposed admissible set is
\begin{equation}
 \mathcal J_i(t_i,\ell)=
 \{j:t_i-\ell\leq e_{ij}\leq t_i,\quad a_{ij}\leq t_i\}.
 \label{exec19:feast-eligible}
\end{equation}
All times use a declared timezone and ordering convention. If feature arrival
and action share a coarse timestamp, an additional sequence number or
conservative ordering is needed to establish which occurred first. The
inclusive inequalities in \eqref{exec19:feast-eligible} are meaningful only
after that convention is fixed. The latest admissible event and revision may
then be selected, provided that the revision order is well defined. A source
row with an unknown availability time does not establish membership in this
set. Conflicting values with the same event and revision coordinates require
an explicit conflict result or a separately justified source ordering.

The distinction is statistical as well as operational. If $\mathcal F_{t_i}$
is the information available to the decision maker, a deployable rule must
be measurable with respect to $\mathcal F_{t_i}$, with any independent policy
randomization stated separately. A retrospectively corrected covariate can
belong to a larger information set. Its association with a later outcome may
be useful for another scientific question, but it cannot silently replace a
predecision predictor, confounder or effect modifier. Nor does successful
retrieval establish exchangeability, positivity, a causal graph or calibrated
uncertainty. Those obligations remain with the analysis that consumes the
retrieved table.

\paragraph{What the selected API and providers actually specify.}
The historical API returns a retrieval job and defaults
\texttt{filter\_by\_created\_timestamp} to false. With the flag enabled,
the requested additional condition is $c_{ij}\leq t_i$ for views that declare
a creation-timestamp column; views without that column are unaffected. The
selected passthrough provider checks the offline store's declared support,
and the base store raises an unsupported-operation error when support is
absent. It would therefore be inaccurate to say that Feast necessarily
ignores creation time. It would also be inaccurate to equate that field with
$a_{ij}$ without evidence about ingestion delays, transformations and the
actual serving path.\citep{exec19_feast_source}

The selected Dask implementation applies the event-time and lifetime filters,
then the optional creation cutoff, then duplicate handling. For late or
unknown creation times it blanks feature-side fields while retaining the
entity-side fields. The accompanying source tests specify the difference
between a later-created value of $0.6$ and an earlier value of $0.4$, exclusion
of null creation times, retention of an entity whose candidates are all late,
and separation of different entities. These are informative upstream test
specifications; they are not native test results obtained by this proposal.
The implementation treats a zero or unspecified lifetime as absence of a
finite lower event-time bound. The proposal must therefore distinguish an
explicitly unbounded history from a positive duration, rather than converting
zero into an instantaneous window.\citep{exec19_feast_source}

The selected BigQuery template applies event-time conditions and, when
requested and configured, the creation cutoff before choosing the latest
event and creation timestamps. Its final join back to feature records uses
those timestamps and the entity-row identifier. If different records share
the chosen coordinates, the proposal cannot infer unique output merely from
the earlier ranking expression. A source uniqueness condition or a qualified
revision key is necessary. The Dask path also needs native checks for equal
timestamps and duplicate decision rows, since sorting and duplicate removal
do not by themselves establish the required population contract. This is a
specific conformance question, not a claim that a defect has been reproduced
in an installed provider.\citep{exec19_feast_source}

\paragraph{An exact decision example with two kinds of late information.}
Consider the following authored finite law, chosen to make the information
constraint explicit. An applicant has an unobserved type $T\in\{g,b\}$ with
equal probabilities. Conditional on action, the potential success indicator
$Y(1)$ has probabilities
\begin{equation}
 P\{Y(1)=1\mid T=g\}=\frac45,\qquad
 P\{Y(1)=1\mid T=b\}=\frac15.
 \label{exec19:feast-law}
\end{equation}
The joint probabilities for $(g,1),(g,0),(b,1),(b,0)$ are respectively
$2/5,1/10,1/10,2/5$. At decision hour $10$, both alternative worlds have the
same observed customer key and the same available score $1/2$; the type label
used to describe the generating law is not an input to the decision. A
successful action has stipulated payoff USD~100, every action costs USD~50,
and no action has value zero. The action value in each type is consequently
\begin{equation}
 v_g=100\frac45-50=30,\qquad
 v_b=100\frac15-50=-30.
 \label{exec19:feast-value}
\end{equation}
These amounts are assumptions of the illustration, not observed bank prices
or estimates from Feast.\citep{exec19_feast_reference}

The original score has event hour $6$, creation hour $7$ and availability hour
$7$. A correction reveals the type-specific score, retains event hour $6$,
and has creation and availability hour $11$. With a four-hour lifetime, the
original lies exactly on the admissible lower event-time boundary. An
event-only retrieval can select the correction despite its later creation;
creation filtering selects the original. In a second history, add a record
with event hour $9$ and creation hour $9$ that reveals the same type-specific
score but is not delivered to the decision system until hour $12$. Now both
event-only and creation-filtered retrieval select that unavailable record.
An actual-availability predicate still selects the original. No timestamp
flag can compensate for a field whose meaning is different from the required
clock.\citep{exec19_feast_reference}

Let the proposed action rule take action precisely when its retrieved score
exceeds $1/2$. The exact comparison is
\begin{center}
\begin{tabular}{p{0.39\linewidth}rrr}
\toprule
History & Event only & Creation cutoff & Availability cutoff\\
\midrule
Correction created after decision & USD~15 & USD~0 & USD~0\\
Additional delivery delay & USD~15 & USD~15 & USD~0\\
\bottomrule
\end{tabular}
\end{center}
The USD~15 entries equal $P(T=g)v_g=(1/2)30$: they evaluate a rule that knows
the type and acts only for $g$. That rule is unavailable at decision hour
$10$. Every feasible rule in this example has the same action probability
$r$ in both worlds, because its available inputs and independent
randomization reveal no type. Hence
\begin{equation}
 E\{\text{net value of a feasible rule}\}
 =r\left(\frac12v_g+\frac12v_b\right)=0.
 \label{exec19:feast-feasible}
\end{equation}
The apparent improvement is thus an information advantage that cannot be
implemented at the declared clock. This derivation does not estimate an
upstream software error or prove a general monetary loss from delayed data.
It supplies a reproducible reason to include availability checks in a
public-data prototype before funding bank experiments.\citep{exec19_feast_reference}

The reference evaluates the entire four-state law with rational arithmetic
and checks invariance to source-row permutations. Its boundary cases separate
the lifetime endpoint, an earlier expired record, a future event, all-late
candidates, unknown clocks, observed nulls, observed zero and conflicting
equal-timestamp values. In an actual cohort, the adapter must also preserve
one result per request identifier. Dropping an applicant because all features
arrive late changes the evaluated population; treating the missing value as
zero changes its meaning. Feature unavailability can itself depend on
customer behavior or servicing workload, so the analysis must address the
resulting selection mechanism explicitly.\citep{exec19_feast_reference}

\paragraph{Historical retrieval, online access and agent tools.}
The selected online API retrieves the latest online feature values, and the
online-store documentation describes materialization into a latest-value
store. A successful online read is therefore not evidence that an old
decision's information has been reconstructed. There is also a separate
provenance obligation for transformations: a feature carrying a predecision
timestamp can still use future observations, a retrospectively fitted
normalizer or a later model version. The proposed library result must bind
the input snapshot, feature and transformation versions, provider, retrieval
mode and relevant clocks. A CLI or MCP wrapper should preserve that same
meaning.\citep{exec19_feast_source}

Feast's selected source distinguishes two MCP surfaces. The standalone
\texttt{feast mcp} process proxies feature and registry servers. Its feature
tools include online retrieval and vector search, but also push and
materialization operations. The selected client turns HTTP error statuses
into tool errors, rather than presenting a denial body as feature data. The
in-process MCP endpoint uses a different extra and deployment path. These
facts justify an integration design, not an assertion that the proposal's
historical contract already exists in either surface.\citep{exec19_feast_source}

For the proposed causal workflow, an agent should request a separately named
historical operation carrying decision identifiers, declared clocks, feature
versions and failure states. Registry browsing can help locate candidate
features; it cannot certify their causal role or availability. Online reads
serve current decisions, while push and materialization change system state.
Those capabilities need separate operational contracts. An unsupported
historical filter, HTTP failure, cancelled job or missing feature must remain
a typed failure or missing observation throughout the library, CLI and MCP
path. Retrying with weaker time semantics would silently answer a different
question.

\paragraph{Reuse choice, alternatives and a bounded public qualification.}
The proposed route is to retain the original Feast implementation behind a
typed historical boundary, initially with one selected provider. The root
source declares Apache-2.0, including redistribution and notice obligations.
The selected project requires Python~3.10 or later and declares thirty direct
dependencies with version ranges; provider extras and the two MCP extras add
different dependencies. The source version is dynamically determined, so a
commit identity alone should not be presented as a verified package-release
number. A complete distribution, resolved environment and applicable
dependency notices remain necessary before executable admission. No native
provider or MCP conformance is established here.\citep{exec19_feast_source}

An explicit bitemporal file or database query is the reference alternative
for a small public prototype: it makes both clocks and the revision rule
visible with fewer components, although its scale and operating costs remain
to be measured. Existing bank feature infrastructure may eventually be the
better integration choice if it already retains the required history and
availability evidence; that cannot be inferred from a feature-store label.
Porting Feast would duplicate provider, timestamp, conversion and failure
semantics without a demonstrated platform requirement. The conditional
preference is therefore integration of an original provider after
qualification, with porting deferred.

Public or generated histories can exercise late corrections, ingestion
delays, null clocks, lifetime boundaries, tied revisions, duplicate requests
and provider errors. The test must compare complete rows and missing states
against a frozen reference, then check that library, CLI and MCP responses
preserve the same retrieval mode. The data owner is responsible for the
meaning of availability and revision fields; the scientific owner for the
information set, population and causal assumptions; and the adapter owner
for conversions, row conservation, failure handling and version changes.
Costs to measure include history retention, joins, validation, data transfer,
materialization and upgrade regressions. These public checks can establish
whether an expensive bank trial is worth proposing. They do not supply bank
truth, measured commercial prices or institutional acceptance.
